\section{Algebraic Protocols, Session Types, and Subtyping}
\label{sec:algebr-sess-types}

This section starts with a brief introduction to algebraic protocols
driven by examples. Next, it motivates constructor subtyping via
compatibility for evolving interfaces. The remaining sections discuss
the translation from TST (traditional, tail-recursive session types)
to algebraic session types with examples, first without subtyping and
finally with subtyping. We use a
syntax inspired by Haskell, as implemented in our artifact. Examples make
liberal use of practical extensions of the formal system (cf.\ \cref{sec:implementation}).

To avoid confusion, we tag each traditional session type name with a
trailing \lstinline+S+ and each protocol type name with a trailing \lstinline+P+.

NOTE: we are going to switch to a structural notation for constructor
subtyping. That is, types can have arbitrary names, but the subtyping
relation is determined by their constructor signatures.

\subsection{Simple Algebraic Protocols}
\label{sec:algebr-datatyp-as}

% Suppose we want to transmit abstract syntax trees for a language with integers
% and addition. Their type might be defined by an algebraic datatype definition as
% in Haskell or ML.
% \begin{lstlisting}
% data Ast = Con Int | Add Ast Ast
% \end{lstlisting}
% This declaration defines the type \lstinline+Ast+ along with the types of the constructor
% functions, which can be used to construct \lstinline+Ast+ values as well as for pattern matching
% on them:
% \begin{lstlisting}
%   Con :        Int -> Ast
%   Add : Ast -> Ast -> Ast
% \end{lstlisting}
% The type \lstinline+Ast+ is recursive and its \lstinline+Add+
% constructor takes two \lstinline+Ast+ values as arguments.
% To define a protocol for transmitting values of type \lstinline+Ast+, we merely change the keyword
% \lstinline+data+ to \lstinline+protocol+ as in
% \begin{lstlisting}
% protocol AstP = ConP Int | AddP AstP AstP
% \end{lstlisting}
% This declaration defines the protocol type \lstinline+AstP+ along with tags \lstinline+ConP+
% and \lstinline+AddP+, which serve to introduce and eliminate branching in the protocol.
% % Tags are encoded by small numbers for transmission.
% The resulting \lstinline+AstP+ protocol can be used in
% two directions, \lstinline+!AstP+ for sending and \lstinline+?AstP+ for receiving
% data.\footnote{The implementation allows us to reuse a
%   \lstinline+data+ declaration as a \lstinline+protocol+ declaration,
%   thus overloading the constructors.} 
% Protocol types only describe behavior, they do not classify run-time values. They inhabit a
% dedicated kind \lstinline+KP+.

% We illustrate the \lstinline+AstP+ protocol with a function to send an
% \lstinline+Ast+ value. We write functions
% in \emph{channel passing style}, which takes a channel as its argument
% and returns it after processing some prefix of the channel's session
% type. The function's type is universally quantified over the continuation
% type \lstinline+s+ of the unexplored part of the channel.
% % This quantification must be introduced explicitly using
% % \lstinline+forall(s:S)+ where \lstinline+KS+ is the kind of session
% % types.
% The argument
% type adds some prefix to \lstinline+s+ that describes the part of the
% session performed by the function. In our example the prefix is \lstinline+!AstP+.
% \begin{lstlisting}
% sendAst : Ast -> forall(s:S). !AstP.s -> s
% sendAst t [s] c = case t of {
%     Con x   -> select ConP [s] c |> sendInt [s] x,
%     Add l r -> select AddP [s] c |> sendAst l [!AstP.s] |> sendAst r [s] }
% \end{lstlisting}

% We use the \lstinline+|>+ operator for reverse function
% application. Defined by \lstinline+x |> f = f x+, the operator associates
% to the left and features low priority, so that \lstinline+x |> f |> g+ means \lstinline+g (f x)+.
% %
% The square brackets indicate type abstraction (on the left) and type
% application (in expressions).
% % By abstracting over a type variable \lstinline+s+ of kind \lstinline+S+
% % (introduced by \lstinline+forall(s:S)+), we require \lstinline+s+
% % to be later replaced by a session type.
% %
% The predefined \lstinline+sendInt+ operation has type
% \lstinline+forall(s:S). Int -> !Int.s -> s+.
% %
% The use of pattern matching with  \lstinline+case+ on algebraic
% datatypes is standard.
% As usual in session type systems, channel values of session type are linear
% and the \lstinline+select+ operation makes an internal choice between alternatives. Selecting a protocol
% constructor sends an encoding of the corresponding tag and ``pushes'' the transmission of the
% fields on the channel type. That is:
% \begin{lstlisting}
%   select ConP [s] : !AstP.s -> !Int.s
%   select AddP [s] : !AstP.s -> !AstP.!AstP.s
% \end{lstlisting}
% A function to process the receiver end of this protocol must take an argument of type \lstinline+?AstP.s+.
% \begin{lstlisting}
% recvAst : forall(s:S). ?AstP.s -> (Ast, s)
% recvAst [s] c = match c with {
%     ConP c -> let (x, c)  = receiveInt [s] c    in (Con x, c),
%     AddP c -> let (tl, c) = recvAst [?AstP.s] c in
%               let (tr, c) = recvAst [s] c       in (Add tl tr, c) }
% \end{lstlisting}
% The \lstinline+match+ on a session-typed channel receives
% a tag and branches according to it. Branching also ``pushes'' transmission of the fields on the
% channel type analogous to the \lstinline+select+ operation:
% \begin{lstlisting}
% c : ?AstP.s |- match c with { ConP c -> ...   -- c : ?Int.s
%                               AddP c -> ... } -- c : ?AstP.?AstP.s
% \end{lstlisting}
% Left of the turnstile $\vdash$, we write the typing of
% \lstinline+c+ before the \matchk. In the body of the \matchk, we write
% the typing of \lstinline+c+ after matching against the selectors \lstinline+ConP+ and \lstinline+AddP+,
% respectively.
% The \lstinline+match c+ (with the type on the left) consumes the linear
% binding for channel \lstinline+c+. The pattern \lstinline+ConP c+
% reintroduces variable \lstinline+c+, binds it with the channel
% after processing the selection and gives it an updated type (on the right) for the body of the
% corresponding branch. The consideration for pattern \lstinline+AddP c+
% is analogous. The choice of the branch is external and depends
% on the received selector tag.

% The protocol in this subsection goes beyond the reach of most session type systems because
% they only support tail recursive session types.
% Thus, functions like \lstinline+sendAst+ and \lstinline+recvAst+ have
% no meaningful type in traditional
% session type systems \cite{DBLP:journals/jfp/GayV10,lindley17:_light_funct_session_types}.
% However, they can be written using context-free or nested session
% types \cite{DBLP:journals/corr/abs-2106-06658,DBLP:conf/icfp/ThiemannV16,DBLP:conf/esop/DasDMP21}. We
% discuss the differences in \cref{sec:related}.


We start with a brief recap of algebraic protocols as presented by \citet{DBLP:journals/pacmpl/MordidoS0V23}.
As a running example, we choose the simple arithmetic server, which is a
standard example from the session type literature (cf.\ 
\cite{DBLP:conf/esop/GayH99}). The arithmetic server
accepts two commands, \lstinline+Neg+ and \lstinline+Quit+. When it
receives the \lstinline+Neg+ command, it waits for an integer
argument, responds with an integer (the result) and terminates. When it receives \lstinline+Quit+, it terminates
the connection directly.

Here is the type of this server adapted from the literature.
\begin{lstlisting}
  Arith0S = &{Neg:?Int.!Int.end, Quit:end}
\end{lstlisting}

The corresponding algebraic protocol \lstinline+Arith0P+ looks as follows.
\begin{lstlisting}
  protocol Arith0P = Quit | Neg Int -Int
\end{lstlisting}
Unlike the session type \lstinline+Arith0S+, \lstinline+Arith0P+ is
just a \emph{protocol}. Here is
the crucial difference: 
The session type \lstinline+Arith0S+ already fixes the direction of communication, as the
operators \lstinline+&+ and \lstinline+?+ indicate receiving. The
protocol \lstinline+Arith0P+ only fixes the structure of the
communication, i.e., \lstinline+Neg+ followed by an \lstinline+Int+
transmission in the same direction while the \lstinline+-Int+ indicates a
transmission in the opposite direction. We materialize the direction
of a protocol by using it in a session type as in
\lstinline+?Arith0P.End+. This algebraic session type corresponds exactly to
\lstinline+Arith0S+. 

The server implementation would match on the command as usual, but
with the following types dictated by the theory of algebraic
protocols.
% \begin{lstlisting}
% c : ?ArithP.s |- match c with { Quit c -> ...  -- c : s
%                                 Neg c -> ... } -- c : ?Int.!Int.?ArithP.s
% \end{lstlisting}
\begin{lstlisting}
c : ?Arith0P.s |- match c with { Quit c -> ...  -- c : s
                                 Neg c -> ... } -- c : ?Int.!Int.s
\end{lstlisting}
The actual implementation of the server is straightforward.
\begin{lstlisting}
arith0Server : ?Arith0P.end -> unit
arith0Server c = match c with {
    Quit c -> close c,
    Neg c  -> let (x, c) = receiveInt [!Int.end] c in
              let c = sendInt [end] (0-x) c in
              close c }
\end{lstlisting}
The type arguments of \lstinline+receiveInt+ and \lstinline+sendInt+
(in square brackets)
are the types of the respective continuation channels. They specify
the type of the variable \lstinline+c+ that is bound in the respective
let.


A client implementation would use the select operator as usual, but
with the following types prescribed by algebraic protocols. This
behavior is dual to the one for matching.

% \begin{lstlisting}
%   select Quit [s] : !ArithP.s -> s
%   select Neg [s]  : !ArithP.s -> !Int.?Int.!ArithP.s
% \end{lstlisting}
\begin{lstlisting}
  select Quit [s] : !Arith0P.s -> s
  select Neg [s]  : !Arith0P.s -> !Int.?Int.s
\end{lstlisting}

A potential client for the protocol would look like this.
\begin{lstlisting}
  arith0Client : !Arith0P.end -> Int -> Int
  arith0Client c x =
    let c = select Neg [end] c in
    let c = sendInt [?Int.end] x c in
    let (r, c) = receiveInt [end] c in
    let _ = close c in
    r
\end{lstlisting}

\subsection{Simple Constructor Subtyping}
\label{sec:simple-constr-subtyp}

Consider an extension of the arithmetic server with an additional
operation indicated by the command \lstinline+Add+. The 
traditional session type changes to include the new alternative.
% \begin{lstlisting}
%   ArithS' = mualpha.&{Add:?Int.?Int.!Int.alpha,Neg:?Int.!Int.alpha, Quit:end}
% \end{lstlisting}
\begin{lstlisting}
  Arith0S' = &{Add:?Int.?Int.!Int.end, Neg:?Int.!Int.end, Quit:end}
\end{lstlisting}
The standard subtyping for session types yields
\lstinline+Arith0S <: Arith0S'+ because a channel of type
\lstinline+Arith0S+ can be processed by code that expects the supertype
\lstinline+Arith0S'+: The code has an additional branch for \lstinline+Add+ that is never
exercised.

The corresponding algebraic protocols are
\begin{lstlisting}
  protocol Arith0P  = Quit | Neg Int -Int
  protocol Arith0P' = Quit | Neg Int -Int | Add Int Int -Int
\end{lstlisting}
With constructor subtyping for algebraic protocols, we can write
\lstinline+Arith0P'[Neg,Quit]+ to select the
constructors \lstinline+Neg+ and \lstinline+Quit+ in
\lstinline+Arith0P'+. 
This way, \lstinline+Arith0P'+ is a supertype of 
\lstinline+Arith0P'[Neg,Quit]+, which corresponds exactly to the
previous protocol \lstinline+Arith0P+ from
\cref{sec:algebr-datatyp-as}.
On the algebraic session type level, we obtain
\lstinline+?Arith0P'[Neg,Quit].end<:?Arith0P'.end+ and, at the client
end,
\lstinline+!Arith0P'.end<:!Arith0P'[Neg,Quit]<:!Arith0P'[Neg]+. The
last subtyping judgment shows that we can reuse our client implementation
with an extended server as \lstinline+arith0Client+ is already typable with
\lstinline+!Arith0P[Neg].end+. 

On the expression level, we extend the match expression in the server
by the case for the \lstinline+Add+ command, just like we would for \lstinline+Arith0S'+.


\subsection{Recursive Algebraic Protocols and Subtyping}
\label{sec:recurs-algebr-prot}

Next, we move to the recursive
version of the arithmetic server. It repeatedly accepts commands
before quitting as indicated by its traditional session type.
\begin{lstlisting}
  ArithS = mualpha.&{Neg:?Int.!Int.alpha, Quit:end}
\end{lstlisting}

The corresponding protocol becomes recursive, too.
\begin{lstlisting}
  protocol ArithP = Quit | Neg Int -Int ArithP
\end{lstlisting}

The implementations of client and  server are omitted because they
change insignificantly with respect to \cref{sec:algebr-datatyp-as}.


% \begin{lstlisting}
% arithServer : ?ArithP.end -> unit
% arithServer c = match c with {
%     Quit c -> close c,
%     Neg c  -> let (x, c) = receiveInt [!Int.?ArithP.end] c in
%               let c = sendInt [?ArithP.end] (0-x) c in
%               arithServer c }
% \end{lstlisting}
% \begin{lstlisting}
%   arithClient : !ArithP.end -> Int -> Int
%   arithClient c x =
%     let c = select Neg [end] c in
%     let c = sendInt [?Int.!ArithP.end] x c in
%     let (r, c) = receiveInt [!ArithP.end] c in
%     let c = select Quit [end] c in
%     let _ = close c in
%     r
% \end{lstlisting}

Instead, we focus on the extension of \lstinline+ArithS+ with
a new recursive case for \lstinline+Add+ as follows.
\begin{lstlisting}
  ArithS' = mualpha.&{Add:?Int.?Int.!Int.alpha,Neg:?Int.!Int.alpha, Quit:end}
\end{lstlisting}
The standard subtyping for recursive session types yields
\lstinline+ArithS <: ArithS'+ because a channel of type
\lstinline+ArithS+ can be processed by code that expects the supertype
\lstinline+ArithS'+: The code has an additional branch for \lstinline+Add+ that is never
exercised.

The corresponding algebraic protocols are
\begin{lstlisting}
protocol ArithP  = Quit | Neg Int -Int ArithP
protocol ArithP' = Quit | Neg Int -Int ArithP' | Add Int Int -Int ArithP'
\end{lstlisting}
With constructor subtyping for algebraic protocols, we can write
\lstinline+ArithP'[Neg,Quit]+ to select the
constructors \lstinline+Neg+ and \lstinline+Quit+ in
\lstinline+ArithP'+. 
This way, \lstinline+ArithP'+ is a supertype of 
\lstinline+ArithP'[Neg,Quit]+, which corresponds exactly to the
previous protocol \lstinline+ArithP+ from \cref{sec:algebr-datatyp-as}. 

On the expression level, we extend the match expression in the server
by the case for the \lstinline+Add+ command, just like we would for \lstinline+ArithS'+.


The subtyping \lstinline+ArithS <: ArithS'+ is decidable by an
$O(n^2)$ algorithm, because the session types involved are traditional
(i.e., tail-recursive). However, the following protocol for serializing abstract
syntax trees is not tail-recursive:
\begin{lstlisting}
  protocol AstP = Const Int | Add AstP AstP | Mul AstP AstP
\end{lstlisting}
Nevertheless, constructor subtyping immediately yields subtypes \lstinline+AstP[Const]+,
\lstinline+AstP[Const,Add]+, and \lstinline+AstP[Const,Mul]+ of
\lstinline+AstP+\footnote{Subtypes without the \lstinline{Const} case
  can also be formed, but they are not very interesting.}. This is remarkable because the non-tail-recursive
nature of the protocol makes it impossible to use the traditional
algorithm to check subtyping of the corresponding (context-free)
recursive types:
\begin{lstlisting}
  AstS1 = mualpha.&{Const: ?Int, Add: alpha;alpha, Mul: alpha;alpha}
  AstS2 = mualpha.&{Const: ?Int, Add: alpha;alpha}
\end{lstlisting}
The current knowledge about subtyping for context-free session types
is that the general problem is undecidable
\cite{DBLP:journals/toplas/Padovani19}, but there exists a
sound semi-decision procedure that seems to work satisfactorily in
practice \cite{DBLP:conf/concur/SilvaMV23}. In our setting, subtyping
is decidable in linear time.

OLD-OLD-OLD

The main feature of algebraic protocols that we did not touch, yet, is
parametric protocols. That means that protocols can be parametric over
other protocols. This feature enables the definition of protocols like
\lstinline+Seq+ from the introduction. We introduce it in XXX along
with the discussion of the translation.


\subsection{From Session Types to Algebraic Protocols}
\label{sec:param-prot}

Starting from  this section we discuss the translation from
traditional, tail-recursive session types to algebraic protocols using
the implementation of \lstinline+arithServer+ in
\cref{sec:algebr-datatyp-as} as a running example.
This code does not use subtyping.

To express all types occurring in \lstinline+arithServer+, we need the
following system of type equations.
\begin{align*}
  \alpha_1 &= \&\{Neg: \alpha_2, Quit: \alpha_3\} & \text{c}\\
  \alpha_2 &= ?\alpha_4.\alpha_5 \\
  \alpha_3 &= \TEnd \\
  \alpha_4 &= \TInt \\
  \alpha_5 &= !\alpha_6.\alpha_1 \\
  \alpha_6 &= \TInt \\
  \alpha_7 &= \alpha_1 \to \alpha_8 & \text{arithServer}\\
  \alpha_8 &= \TUnit \\
  \alpha_9 &= \alpha_3 \to \alpha_8 & \text{close} \\
  \alpha_{10} &= \alpha_2 \to \alpha_{11} & \text{receiveInt}\\
  \alpha_{11} &= \alpha_4 \times \alpha_5 \\
  \alpha_{12} &= \alpha_5 \to \alpha_1 \\
  \alpha_{13} &= \alpha_6 \to \alpha_{12} & \text{sendInt}
\end{align*}

For this system, typing rules have the form $\mathcal{E}\mid \Gamma
\vdash e : \alpha$ where $\mathcal{E}$ is a system of flat type
equations, $\Gamma$ is a typing context mapping variables to type
variables defined in $\mathcal{E}$, $e$ is an expression, and $\alpha$
is a type variable. Here are some exemplary typing rules for variables,
function application, and abstraction:
\begin{mathpar}
  \inferrule{}{\mathcal{E} \mid \Gamma, x:\alpha \vdash x : \alpha}

  \inferrule{
    \mathcal{E} \mid \Gamma \vdash \termc{e_1} : \alpha_1 \\
    \mathcal{E} \mid \Gamma \vdash \termc{e_2} : \alpha_2 \\
    \alpha_1 = \alpha_2 \to \alpha_3 \in \mathcal{E}
  }{
    \mathcal{E} \mid \Gamma \vdash \appe{e_1}{e_2} : \alpha_3
  }

  \inferrule{
    \mathcal{E} \mid \Gamma, x:\alpha_2 \vdash \termc{e} : \alpha_3 \\
    \alpha_1 = \alpha_2 \to \alpha_3 \in \mathcal{E} \\
    \typec{T} \equiv \emph{extract}(\mathcal{E}, \alpha_2)
  }{
    \mathcal{E} \mid \Gamma \vdash \abse{x}{T}{e} : \alpha_1
  }
\end{mathpar}
The typing rule for typed lambda abstraction has to extract the type
from the equation system by traversing it. For each session type
variable $\alpha_i$, extraction inserts a $\mu\alpha_i$ binder and
refers back to $\alpha_i$ if it is seen again on the traversal.

Any recursively typed term can be type checked using a system of
recursive type equations, so nothing is lost in this step.

In our running example, minimization equates variables $\alpha_4$ and
$\alpha_6$, both of which stand for the type $\TInt$. No further
simplification is possible.

Next, we focus on the session type variables $\alpha_1$, $\alpha_2$,
$\alpha_3$, and $\alpha_5$. We start with $\alpha_1$ because it holds
a branch type and invent a new protocol name, say,
\lstinline+NegQuit+ with two constructors \lstinline+Neg+ and
\lstinline+Quit+.  For the \lstinline+Neg+ constructor, we continue
with $\alpha_2 = ?\alpha_4.\alpha_5$. As this session type has the
same direction as branching, we extract its payload type $\alpha_4$
(i.e., $\TInt$) and add it to the constructor with positive
direction. Continuing with $\alpha_5 = !\alpha_6.\alpha1$, we find a
session type in the reverse direction, so we extract its payload from
$\alpha_6$ (i.e., $\TInt$) and add it to the constructor with negative
direction. Continuing with $\alpha_1$, we realize that we have seen
this type variable before and conclude the constuctor with a positive
\lstinline+NegQuit+ argument.

For the \lstinline+Quit+ constructor, we consider $\alpha_3 =
\TEnd$. However, we cannot add $\TEnd$ as a protocol argument because
it would mean that we wanted to transmit a \emph{channel} of type $
TEnd$, so we do not add it as an argument to the constructor, but
rather finish the constructor without arguments.

The resulting protocol type is
\begin{lstlisting}
  protocol NegQuit = Neg Int -Int NegQuit | Quit
\end{lstlisting}

\subsection{Discussion}
\label{sec:constr-subtyp-algebr}

Consider an extension of the arithmetic server with an additional
operation indicated by the command \lstinline+Add+.
Here is an example that demonstrates the limitations of this approach
to protocol subtyping.
\begin{lstlisting}
  protocol Msg = M1 | M2
  protocol QR a b = Step a -b (QR a b) | Quit
\end{lstlisting}
The protocol \lstinline+QR Msg Msg+ indicates that client and server
send \lstinline+M1+ and \lstinline+M2+ messages forth and back. If we
evolve the message protocol with additional messages, then constructor
subtyping will fail. Consider the extended message type
\begin{lstlisting}
  protocol Msg' = M1 | M2 | M3
\end{lstlisting}
Using \lstinline+QR Msg' Msg'+ means that the client can send an
arbitrary message \lstinline+M1+, \lstinline+M2+, or \lstinline+M3+,
but it has to be prepared to handle the same messages as input. That
is, even a client that promises to send only \lstinline+M1+ and
\lstinline+M2+ has to be prepared to process a message \lstinline+M3+
in response. Hence, \lstinline+?QR Msg'[M1,M2] Msg'+ is a subtype of
\lstinline+?QR Msg' Msg'+, but \lstinline+?QR Msg'[M1,M2] Msg'[M1,M2]+
is not, because the second parameter of \lstinline+QR+ occurs under
inversion (i.e., with revered direction). In short, we cannot obtain a
type equivalent to \lstinline+QR Msg Msg+ as a subtype of
\lstinline+QR Msg' Msg'+.

\begin{lstlisting}
  protocol AB = A | B
  protocol Msg'' a b = M1 -a | M2 -b
  protocol Msg''' a b c = M1 -a | M2 -b | M3 -c

  !Msg'' AB A <: !Msg'' AB AB
  ?Msg'' AB A <: ?Msg'' A A
\end{lstlisting}

%%% Local Variables:
%%% mode: latex
%%% TeX-master: "main"
%%% End:
